\hypertarget{scenario-1-building-a-model-of-cancer-cell-fate-in-response-to-tnf-from-scratch}{%
\section{Scenario 1: Building a model of cancer cell fate in response to
TNF from
scratch}\label{scenario-1-building-a-model-of-cancer-cell-fate-in-response-to-tnf-from-scratch}}

\hypertarget{prompt}{%
\subsection{Prompt:}\label{prompt}}

You are my multiscale modelling assistant. Together we are creating a
mutliscale multicellular model of cancer cell fate in response to TNF.
In order to chieve this we are going to use NeKo, MaBoSS and PhysiCell
to create a simulation where an initial group of cancer cell will grow
while TNF diffuses in the microenvironment. When the cells sense the TNF
they will activate the input nodes of their boolean network. According
to the final state of the network, they will perform behaviors such as
die for apoptosis or proliferte. Those are he steps to follow: 1) create
a network using Omnipath that describes cancer cell fates in response to
tnf. This network should contains TNF sensing receptr as well asgenes
responsible for proliferation and apoptosis. Once create the network we
have to remove the bimodal interactions and check the connectivity. If
the network is connected, export it in bnet file. After that generate
the maboss file (.bnd and .cfg), create a simulation and select some
nodes to identify the output, to represent apoptotic state and
proliferative state. Once set those nodes as output, we run the maboss
simulation, get the results and study the outcome (does the network make
sense? are the states biologically feasible?). This is the first part.
In the second part we will build the PhysiCell simulation and add maboss
to the cancer cell type. Is everything clear?

\hypertarget{final-results}{%
\subsection{Final results:}\label{final-results}}

\textbf{Multiscale Simulation Summary} Session ID: f6c3f9c5\ldots{}
Scenario: Tumor spheroid in a TNF‐diffusive microenvironment. Cancer
cells sense TNF via TNFRSF1A, run a MaBoSS Boolean network, and execute
apoptosis or proliferation based on CASP3 and CCND1 outputs.

---- Network construction ---- • Seed genes: TNF, TNFRSF1A, FADD, RIPK1,
TRAF2, NFKB1, RELA, BCL2, BAX, CASP3, CASP8, CCND1, MYC • Nodes: 53
Interactions: 332 • Removed 14 bimodal interactions, confirmed full
connectivity • Exported BNET: \texttt{Network\_1.bnet}

---- MaBoSS model ---- • Generated \texttt{output.bnd} +
\texttt{output.cfg} • Simulation parameters: time\_tick=0.5,
max\_time=1000, samples=10000 • Output nodes: CCND1 (proliferation),
CASP3 (apoptosis)

Simulation results (final state probabilities):

\begin{enumerate}
\def\labelenumi{\arabic{enumi}.}
\tightlist
\item
  \textbf{Baseline} • CCND1-ON = 1.00 • CASP3-ON = 0.00
\item
  \textbf{Mutation BCL2 = OFF} • CCND1-ON = 1.00
\item
  \textbf{Mutation CASP8 = ON} • CCND1-ON = 1.00
\item
  \textbf{Mutation BAX = ON} • CCND1-ON = 1.00
\end{enumerate}

\emph{All runs ended in the proliferative attractor.}

---- PhysiCell integration ---- • Domain: 1000×1000×1000 μm, mesh 20 μm,
duration 7200 min • Substrate added: TNF (D=1000 μm²/min, decay=0.1) •
Cell type: \texttt{cancer\_cell} (Ki67\_basic) • PhysiBoSS: --
Intracellular model attached (\texttt{output.bnd}/\texttt{.cfg}) -- dt=6
min, scaling=1, stochasticity=0 -- Input link: intracellular TNF → TNF
activation -- Output links: • CASP3 → apoptosis • CCND1 → cycle entry •
Cell rules mirror Boolean outputs: -- CASP3↑→ apoptosis; CASP3↓→ cycle
entry -- CCND1↑→ cycle entry; CCND1↓→ apoptosis

\hypertarget{scenario-2-building-a-set-of-rules-for-an-existing-physicell-simulation-by-extracting-the-context-directly-from-the-configuration-file}{%
\section{Scenario 2: Building a set of rules for an existing PhysiCell
simulation by extracting the context directly from the configuration
file}\label{scenario-2-building-a-set-of-rules-for-an-existing-physicell-simulation-by-extracting-the-context-directly-from-the-configuration-file}}

\hypertarget{prompt-1}{%
\subsection{Prompt:}\label{prompt-1}}

I have a PhysiCell configuration file that I would like to extend with
biologically accurate rules. Load the PhysiCell configuration file,
extract the biological scenario, extract the substrates, cell types and
custom parameters that compose the simulation and, carefully examining
the available signals and behaviors, add a series of rules to improve
the model.

\hypertarget{final-results-1}{%
\subsection{Final Results:}\label{final-results-1}}

I've integrated biologically informed cell‐behavior rules into your
configuration and exported the updated XML as
\texttt{PhysiCell\_settings\_with\_rules.xml}. It includes:

\begin{itemize}
\tightlist
\item
  Tumor hypoxia‐induced necrosis and apoptosis
\item
  Macrophage migration toward debris, phagocytosis, and differentiation
  (M0→M1/M2)
\item
  M1/M2 secretion loops of pro-/anti-inflammatory factors
\item
  Effector T-cell tumor attack on contact
\end{itemize}

Next, place the new XML in your PhysiCell project and run your
simulation engine with it, for example:

\begin{Shaded}
\begin{Highlighting}[]
\ExtensionTok{./myproject}\NormalTok{ PhysiCell\_settings\_with\_rules.xml}
\end{Highlighting}
\end{Shaded}

\hypertarget{scenario-3-iteratevely-refine-a-maboss-model-through-neko-using-the-basic-knowledge-of-the-llm}{%
\section{Scenario 3: iteratevely refine a MaBoSS model through NeKo
using the basic knowledge of the
LLM}\label{scenario-3-iteratevely-refine-a-maboss-model-through-neko-using-the-basic-knowledge-of-the-llm}}

\hypertarget{prompt-2}{%
\subsection{Prompt:}\label{prompt-2}}

I want to create a Boolean model that describes how TNF affects cancer
cell fate. In order to do so, I want you to create a network using
Omnipath as interactions database. Once created the network, remove the
bimodal interactions, export the network in bnet format and convert it
into a MaBoSS model. Run the MaBoSS model and het the results. If the
results are biologically accurate and relevant, test some known
mutations end verify the accuracy of the results. If the network is not
satisfying, check important pathways within the network and individuate
nodes or interactions that may not be relevant and bring inaccuracies in
the model. Finally, make sure the final network always contains relevant
nodes, important for our scenario. Iterate the process until the MaBoSS
model is accurate enough.

\hypertarget{final-results-2}{%
\subsection{Final results:}\label{final-results-2}}

Here's a concise summary of our workflow and key findings:

\begin{enumerate}
\def\labelenumi{\arabic{enumi}.}
\tightlist
\item
  Network construction• Seed genes: TNF, TNFRSF1A, TRADD, FADD, CASP8,
  CASP3, BCL2, NFKB1 (+ later TNFRSF10B, CFLAR, PTEN)• Database:
  Omnipath, only signed interactions, 2-step paths (later iterations
  used bias/consensus settings)• Pruned all ``bimodal'' edges and
  low-impact nodes for clarity
\item
  Model translation \& setup• Exported to \texttt{.bnet}, converted to
  MaBoSS \texttt{.bnd}/\texttt{.cfg}• Thread count set to 10, 10 000
  samples, max\_time=1000, time\_tick=0.5
\item
  Baseline simulation (TNF input)• Unique attractor: full activation of
  NFκB‐driven survival nodes, no apoptosis (CASP3) or necroptosis (MLKL)
\item
  Single-gene perturbations• BCL2 OFF, NFKB1 OFF, CASP8 OFF, RIPK1 OFF,
  CFLAR OFF, XIAP OFF → no switch to apoptotic or necroptotic attractors
\item
  Co-stimulation \& inhibitor scenarios• TNF + TRAIL (TNFRSF10B ON) →
  same survival attractor• XIAP OFF + TRAIL → still survival‐dominant
\item
  Added PTEN• No impact on attractor
\item
  Interpretation • The network is robust: single KOs or single-ligand
  changes aren't enough to trigger death • Likely need combined
  perturbations (e.g., CFLAR+XIAP double KO) or inclusion of additional
  regulators (e.g., CFLAR, CASP3 co-activators)
\end{enumerate}

Next steps could include systematic double-mutations or quantitative
comparison of CASP3 activation probabilities against experimental values
(≈0.6--0.8 under TNF+TRAIL). Let me know which scenario you'd like to
explore. GitHub Copilot
